\scanpage{25}
\mspage{12}

qui fait une restriction du groupe structural à $V_J(\mathbb F_2)$, d'où $\bigwedge_{i\in J}D_i\simeq\mathord{1\mkern-4mu\rule{.055em}{.66em}}$, OK.

Il y aurait lieu aussi d'interpréter les $10$ classes d'isomorphie d'espaces homogènes $\ne(1)$ sous $G\simeq\operatorname{Aut}(\text{Cube})$ (correspondant à des \deleted{\unsure{deux|des}} sous-groupe d'isotropie $\ne G$ \unsure{\}|)}, définis à conjugaison près) en termes d'\deleted{objets} éléments de structure associés au cube (tels sommets, arêtes, faces, ou \deleted{tels} objets mod antipodisme -- i.e. sur le cube gauche -- ou \deleted{\unsure{couples de points|repères de points}} \inserted{repères} directs -- ce qui en fait déjà sept \inserted{\unsure{non iso|en iso.}}, de cardinaux $8,12,6,4,6,3,24$. Il y en a un de cardinal $2$, quotient de l'ens des sommets (correspondant à l'ens. $\underline\omega_I$), un autre d'ordre $12$ correspondant aux \deleted{\unsure{fa}} diagonales des faces \inserted{(ou codiagonales, si on préfère)} (les stabilisateurs \deleted{\unsure{de ces él. d'ordre}} \inserted{correspondent} aux él. d'ordre $2$ de $G$ qui \deleted{\unsure{ne}} sont \deleted{\unsure{pas}} dans $V$), enfin un troisième qui est formé des paires d'éléments distincts de $J$ (i.e. repères de $J$), $J$ étant interprété comme l'ens des ``directions de \deleted{\unsure{plans}} \inserted{faces}'' (\deleted{\unsure{ou ens}} (ens de faces/antipodisme) ou des directions d'arêtes. Le compte y est !

% slanted
\margin{position relative de deux sommets à dist. comb. $1$, ou encore de deux arêtes : dist. comb. $1$, \deleted{ou à dist. comb $2$,}}

\deleted{\unsure{5|En}}

% archive sheet number: 24

